% !TeX root = ../../surface-lagrangian-torus-fibration.tex
% 2026-11-21 / v1.3: 平均順序の相対化，normal flow の自然性と全域接着への入力．
\section{複数の三重点 end における平均順序の相対化}
\label{ke:relative-corner-section}

三つの end で異なる primitive normal subgroup を先に平均したい場合，
一つの subgroup を指定する系\ref{ke:staged-closed-corner-gh}だけでは足りない．
本節では，互いに離れた end 上の二段階構成を，
一つの actual Lagrangian torus bundle に同時に組み込む．
そのために命題\ref{ke:staged-hamiltonian-comparison}の比較 Hamiltonian 自体を評価する．
座標の滑らかな補間によって作用が非正則になる場合も，
同じ toric 模型に対する小摂動という仮定の下で扱う．
Neck 全体との接着に必要な moving domain と共通 moment interval は，
本節の相対化だけからは得られず，
定理\ref{ke:neck-corner-diagonal-matching}で exhaustion とともに選ぶ．

\subsection{小摂動の平均と比較 Hamiltonian}

$U$ に自由な $T^2$ 作用を与え，その全平均を $\mathcal A$ とする．
$\omega_0=d\lambda_0$ を作用不変な K\"ahler 模型とし，
$h=g_{\omega_0}$ と置く．$h$ と $\lambda_0$ も作用不変とする．
従って作用は模型の複素構造 $J_0$ を保存する．
以下の微分の評価は，角度の universal cover 上の固定半径正規化 chart で用いる．
$h,h^{-1}$ の固定階の bounds は一様とし，torus quotient の injectivity radius の
一様下界は要求しない．有限個の固定整数基底だけを用いる．

\begin{proposition}[比較 Hamiltonian の定量的な primitive 評価]
\label{ke:quantitative-order-hamiltonian}
$\omega=d\lambda$ を閉 $2$-form とし，ある $0<e\le1$，十分小さい $\delta>0$ に対し
\begin{equation}\label{ke:eq:order-primitive-input}
 \|\lambda_0\|_{C^4(h)}\le C_0,\qquad
 \|\omega-\omega_0\|_{C^4(h)}+
 \|\lambda-\lambda_0\|_{C^4(h)}\le\delta,\qquad
 \|V_j\|_{C^4(h)}\le C_0 e
\end{equation}
を仮定する．$V_j$ は用いる有限個の周期 $1$ の生成場である．
すべての平均の chain homotopy に centered kernel $1/2-s$ を用いる．

一つの整数基底の一段階 Moser map を $\Phi$ とし，別に指定した整数基底での
normal-first 二段階 map を $\Theta$ とする．
内外の領域には共通の $h$-buffer $b>0$ を取り，$Ce\delta<b$ とする．
これらの領域は torus 飽和したものとする．
必要なら内側領域を一度縮めれば，$F_0=\Phi$，$F_1=\Theta$ である
endpoint maps の滑らかな族 $F_r$ を構成でき，
\[
 F_r^*\omega=\mathcal A\omega,\qquad
 F_r^*\lambda=\mathcal A\lambda+dQ_r
\]
を満たす．共通の像領域上で $\Psi_r=F_r\Phi^{-1}$ は $\omega$-Hamiltonian であり，
その生成場 $W_r$ と，$\iota_{W_r}\omega=dH_r$ の規約で選んだ Hamiltonian は
\begin{equation}\label{ke:eq:order-hamiltonian-bound}
 \|W_r\|_{C^1(h)}+\|H_r\|_{C^2(h)}\le C e\delta
\end{equation}
を満たす．定数は模型の一様幾何と固定整数基底に依存するが，領域の直径には依存しない．
作用が別途与えられた actual 複素構造を保存すること，および
$\omega$ が模型の複素構造 $J_0$ に関して型 $(1,1)$ であることは不要である．
\end{proposition}

\begin{proof}
各 $K_j$ は centered homotopy，$\mathcal A_j$ は円平均とする．
模型とその primitive の不変性により
\[
 K_j\omega_0=0,\qquad K_j\lambda_0=0.
\]
$h$ に関する平均は固定階の tensor norms を増やさず，contraction の微分には
$V_j$ の同じ階以下の微分だけが現れる．従って
\[
 \|K_j\omega\|_{C^3(h)}+
 \|K_j\lambda\|_{C^3(h)}\le Ce\delta.
\]
先に他の円平均を施した場合にも同じ評価が成り立つ．
各平均とその凸結合は $\omega_0$ から $C\delta$ 以内にある．
特に模型の複素構造 $J_0$ に対して
\[
 \Omega(v,J_0v)\ge(1-C\delta)|v|_h^2
\]
を満たし，閉形式 $\Omega$ は非退化である．この箇所では平均の K\"ahler 性を使わない．

指定した基底について $\eta=\mathcal A_1\omega$，$\kappa=\mathcal A\omega$ とし，
命題\ref{ke:staged-hamiltonian-comparison}の三角形上の forms
\[
 \Omega_{u,v}=(1-u)\kappa+(u-v)\eta+v\omega,\qquad
 \Lambda_{u,v}=(1-u)\mathcal A\lambda+(u-v)\mathcal A_1\lambda+v\lambda
\]
を用いる．$0\le v\le u\le1$ の中で，対角 path と二辺を通る path を
滑らかな族 $(u_r(s),v_r(s))$ で結ぶ．二辺の接続を stationary に parameter 化すれば，
必要な $r,s$ 微分を一様に有界にできる．
\[
 B_{r,s}=\dot u_r K_2\mathcal A_1\omega+\dot v_r K_1\omega,
 \qquad \iota_{Y_{r,s}}\Omega_{u_r,v_r}=B_{r,s}
\]
と置く．$\Omega$ の逆の bounds と上の contraction 評価から，
$Y_{r,s}$，$\partial_rY_{r,s}$ の $C^2(h)$ norms は $Ce\delta$ 以下である．
この長さと共通 buffer によって全 flow が存在する．
Flow の変分方程式と Gronwall の不等式から，endpoint の変位，微分の恒等写像からの差，
および $\partial_rF_r$ の $C^2$ norm も同じ率で制御できる．
Backward flow にも同じ長さ評価を用いれば，共通の内側集合を全 $F_r$ の像に含められる．

\[
 C_{r,s}=\dot u_rK_2\mathcal A_1\lambda+\dot v_rK_1\lambda
\]
とすると $B_{r,s}=-\partial_s\Lambda_{u_r,v_r}-dC_{r,s}$ である．
時刻 $s$ までの flow を $F_{r,s}$ と書けば，exact transport の式は
\[
 Q_r=\int_0^1F_{r,s}^*
       \bigl(\Lambda_{u_r,v_r}(Y_{r,s})-C_{r,s}\bigr)\,ds
\]
で与えられる．$\Lambda$ の norm は一様に有界で，
$C_{r,s}$ とその $r$ 微分は $C^2(h)$ で $O(e\delta)$ である．
また $\partial_r\Lambda$ は $O(\delta)$ である．模型部分の係数の和が常に $1$ なので，
この微分で大きな模型項は現れない．従って
\[
 \|Q_r\|_{C^2(h)}+\|\partial_rQ_r\|_{C^2(h)}\le Ce\delta.
\]
Exact identity を $r$ で微分し，$W_r=(\partial_rF_r)\circ F_r^{-1}$ と置くと
\[
 F_r^*(\iota_{W_r}\omega)
 =d\bigl(\partial_rQ_r-F_r^*(\lambda(W_r))\bigr).
\]
従って
\[
 H_r=(\partial_rQ_r)\circ F_r^{-1}-\lambda(W_r)
\]
が求める Hamiltonian であり，\eqref{ke:eq:order-hamiltonian-bound}を得る．
単に $dH_r$ を評価して基点から積分する方法を用いないため，領域の直径は現れない．

最初の整数基底が指定基底と異なる場合には，一段階の二つの homotopies $H,H'$ に対し
同じ線形 form path 上で primitive を $(1-r)H\omega+rH'\omega$ と補間する．
$H\omega=\mathcal A\lambda-\lambda-d(H\lambda)$ とその $H'$ 版により，
上と同じ exact transport と評価が成立する．この族と三角形の族を，
parameter の端で stationary にしてつなげれば一般の場合を得る．
\end{proof}

\begin{theorem}[互いに素な end で指定した二段階構成の相対統合]
\label{ke:relative-corner-patch}
命題\ref{ke:quantitative-order-hamiltonian}の仮定の下で，有限個の end ごとに
primitive normal subgroup とその整数基底への補完を指定する．
各 end の torus 飽和した内側領域を $E_a^-$ とし，少し広い領域 $E_a^0$，$E_a^+$ を取る．
$\overline{E_a^-}\subset E_a^0\Subset E_a^+$ とし，$E_a^+$ は互いに素で，
すべての比較 flow と inverse compositions の共通領域に相対コンパクトに含まれるとする．
特に $\overline{E_a^+}$ は全ての $F_{a,r}$ の共通の像の内部に含める．
これらを同じ ambient $U$ の部分集合として取り，source の $E_a^-$ と
actual 像を置く $E_a^0$ の間には少なくとも $2Ce\delta$ の余裕を置く．
$\chi_a=1$ on $E_a^0$，$\operatorname{supp}\chi_a\Subset E_a^+$，
$\|\chi_a\|_{C^3(h)}\le B$ なる cutoffs を取る．
入れ子領域間の $h$-buffer が $C_B e\delta$ より大きければ，
一段階構成の torus bundle を actual $\omega$-Hamiltonian isotopy で変形し，
各 $E_a^-$ の torus 上で，指定した normal-first map $\Theta_a$ の像と
marked cycles ごと一致する一つの Lagrangian bundle を得る．
補正は $\bigcup_a E_a^+$ の外で恒等であり，変位と微分の変化は $O_B(e\delta)$ である．
\end{theorem}

\begin{proof}
各 $a$ に対して前命題を用い，$\Psi_{a,r}=F_{a,r}\Phi^{-1}$ と $H_{a,r}$ を取る．
共通の像領域で
\[
 \widetilde H_r=\sum_a\chi_a H_{a,r}
\]
を定め，support の外へ零で延長する．
$\iota_{X_{\chi H}}\omega=d(\chi H)$ より
\[
 X_{\chi H}=\chi X_H+H X_\chi,
\]
従ってその $C^1$ norm は $C_B e\delta$ 以下である．
この compact support を持つ時間依存 Hamiltonian の flow を $\Upsilon_r$ とする．
Buffer と前命題の変位評価により，$\Phi(E_a^-)$ から始まる
$\Psi_{a,r}$ の全 trajectories は $E_a^0$ に留まる．
そこで $\chi_a=1$，他の cutoffs は零なので，ODE の一意性により
\[
 \Upsilon_1\Phi(x)=F_{a,1}(x)=\Theta_a(x)\qquad(x\in E_a^-).
\]
$\kappa=d(\mathcal A\lambda)$ と primitive の不変性から
$\kappa(V_1,V_2)=0$ であり，$\Phi$ の全 torus は $\omega$-Lagrangian である．
$\Upsilon_1$ は $\omega$-symplectic なので，その像は一つの proper torus bundle である．
元の torus を同じ写像で移しているため，表示された一致は marked cycles も保つ．
Hamiltonian に共通の primitive が定義されている領域では，
$\Upsilon_1^*\lambda-\lambda$ は exact であり，全 cycle periods を保つ．
最後の変位と微分の評価は $X_{\widetilde H_r}$ の bounds から従う．
\end{proof}

\subsection{Actual triple での入力と適用範囲}

本節の primitive の仮定は，metric の小摂動だけから任意の primitive について
結論するものではない．Actual triple では選択を次のように固定できる．
$\Omega=d\log\zeta_1\wedge d\log\zeta_2$，
$\omega_0=k_L$，$h=g_L^{\mathrm{sf}}$ とし，
\[
 \lambda=d^c\log\frac{\omega_t^2}{|\Omega|^2},\qquad
 \lambda_0=d^c\log\frac{k_L^2}{|\Omega|^2}
\]
を取る．関連 KE 主稿\cite[\texttt{coarse:thm:full-triple-tensor}]{KELocalCurrent}により，
$\min\ell_i\ge m$ の少し広い領域で全固定階の metric 誤差は $O(m^{-1})$ である．
同じ正規化 chart の determinant を取れば
$v_t=\log(\omega_t^2/k_L^2)$ の全固定階の norm も $O(m^{-1})$ となり，
\[
 \lambda-\lambda_0=d^cv_t=O_{C^4(h)}(m^{-1})
\]
を得る．模型 primitive は $\lambda_0=2d^c\phi_L$ である．
同稿 \texttt{coarse:lem:simplex-model} の potential bounds により，
固定相対窓で $\phi_L$ の振幅は一様に有界である．凸性で一階微分を評価し，
模型 metric の一様な高階 bounds を合わせれば $\lambda_0$ の $C^4$ bound を得る．
また固定整数生成場の $C^4$ norm は $O(m^{-1})$ である．
従って $e=C/m$，$\delta=C/m$ として本節を適用でき，比較 Hamiltonian は
$O(m^{-2})$ となる．

さらに命題\ref{ke:three-arm-coordinate-interpolation}の
逆向きの座標同定 $G=F_1^{-1}$ を用いる場合も，
同じ正規化 chart で $G-\mathrm{id}=O_{C^5}(\epsilon)$ なら，
$G^*\omega_t$，$G^*\lambda$ に対して
$\delta=C(m^{-1}+\epsilon)$ が使える．作用が非正則になる補間帯では，
正値性は本節の共通模型による tame 評価から従う．
$G^*\lambda_0-\lambda_0$ の $C^4$ 評価には模型 primitive の $C^5$ bound を用いる．
上の模型の議論は任意の固定階に適用できるので，この追加の一階も一様である．
三つの dominant arms の支持を固定の深度比で離して選べば，
深度比の対数の滑らかな関数により一様な $C^3$ cutoffs を取れる．
したがって $m=R$ の closed-corner 構成では，
各 arm に指定した normal-first 構成を一つの bundle へ組み込める．
補正の変位と $C^1$ 誤差は $O(R^{-2}+\epsilon/R)$ であり，
$\epsilon$ が指数的に小さい場合，既存の内在・protected GH 誤差
$O(\log R/R)$ を変えない．
実際，同じ softmax domain の射影を用い，補正した parametrization で actual metric を
引き戻すと，模型との差は引き続き $O(R^{-1})$ である．内在距離の上界には
同じ模型の square lifts を用いる．Protected 距離の下界では外側領域全体の actual/model
比較は変わらず，補正した端点からの clipped-depth 誤差は
$O(R^{-3}+\epsilon/R)$ である．必要な source と像の余裕は，
$\min\ell_i\ge R$ と $\min\ell_i>R/2$ の間に取れる．
従って定理\ref{ke:closed-corner-gh}と
系\ref{ke:closed-corner-protected-gh}の証明を同じ率で適用できる．
これは同一 triple の局所結果である．

一般の $m_t\to\infty$，$\log m_t=o(R_t)$ に対しても，
外側境界の cutoffs を $\log(\ell_a/m_t)$ の固定幅で取り，
定理\ref{ke:slow-cut-closed-corner-gh}の同じ softmax domain に適用できる．
このとき座標補間を含めた pullback metric の誤差は $O(m_t^{-1})$，
平均順序の補正は $O(m_t^{-2})$ である．
Protected 距離の端点誤差には同定理の pruning 定数 $O(R_t^{-1})$ を掛ければよい．
従って三端の指定を組み込んだ一つの bundle にも
\eqref{ke:eq:slow-cut-distortion}と同じ局所 distortion 評価が成立する．

系\ref{ke:transition-normal-primitive}により，内側 arm では射影的 neck と
triple の canonical primitive の差は，残余基底のみに依存する閉形式と
一価関数の完全微分の和になる．従って円平均後の normal moment は厳密に一致する．
第二の cycle の period まで零差になるとは限らない．
全 marked periods を比較するときは実際に同じ primitive を用いる必要がある．
本節の小さい Hamiltonian の評価には上記の小さい triple primitive を用い，
neck primitive への置換後にも同じ norm bound があるとは仮定しない．
また各 end の正則作用が projected neck の作用と一致していても，
normal trajectories が両側の共通開領域に留まることを確認する必要がある．
固定残余基底の neck から三重点へ動く collar までの被覆と，
source-inclusive neck に共通の moment interval・開領域は
定理\ref{ke:neck-corner-diagonal-matching}で確保し，
系\ref{ke:neck-corner-saturated-union}で全交わりの飽和性を確認する．
全曲面の距離比較は，後続の定理\ref{ke:global-ambient-lower}，
\ref{ke:global-ambient-upper}による．
単一 arm の interval と reduced 葉については，以下で実際の評価を与える．

\subsection{Reduced collar の graph による延長}

\begin{lemma}[小さい circle graph の相対延長]
\label{ke:reduced-collar-graph}
境界付き曲面 $P$ の互いに素な collars に座標 $(a,\theta)\in J\times S^1$ を取り，
基準 foliation を $a=\mathrm{const}$ とする．
共通 moment parameter $\mu\in I$ ごとに指定された end の葉が
\[
 a=Q_\mu(c,\theta),\qquad
 \|Q_\mu(c,\theta)-c\|_{C^1(c,\theta)}\le\varepsilon
\]
という graph で，$\mu$ に滑らかに依存するとする．
開 end germ 上で $1$ となり，少し広い collar に support を持つ cutoff $\chi(c)$ があり，
$\varepsilon(1+\|\chi'\|_\infty)<1/2$ とする．
Graph は cutoff の support を含む広い collar 全体で定義されているものとする．
このとき基準の radial 葉を指定 end germs へ写す $\Diff_0(P)$ の滑らかな族が存在し，
広い collars の外では恒等である．特にその外の sources と saddle markings を固定できる．
\end{lemma}
\begin{proof}
各 collar で
\[
 D_\mu(c,\theta)
 =\bigl(c+\chi(c)(Q_\mu(c,\theta)-c),\theta\bigr)
\]
と定め，外へ恒等写像で延長する．第一成分の $c$ 微分は $1/2$ より大きい．
また collar の両端近くで恒等であるため，全体は orientation を保つ diffeomorphism である．
$\chi$ を $s\chi$，$0\le s\le1$ に置き換えれば，支持を保つ isotopy を得る．
互いに素な collars で同時に行え，parameter $\mu$ にも滑らかである．
$\chi=1$ の開 collar 上では基準円の像が指定した graph 葉になる．
\end{proof}

この補題を系\ref{ke:staged-reduced-extension}へ使うには，actual reduction の共通同定と，
その座標での graph 表示・$C^1$ 評価を確認しなければならない．
Ambient Moser map が小さいという事実だけを，未指定の reduced 座標での
一様な graph 評価に読み替えない．一次元 reduced 葉の Lagrangian lift を作ることと，
第二の action や critical actions をあらかじめ指定した値へ合わせることも区別する．

\begin{proposition}[単一 dominant arm の actual reduced graph]
\label{ke:single-arm-reduced-graph}
一つの source-free arm の標準積座標を $\zeta_i\zeta_j\zeta_k=t$ とし，
\[
 a=\ell_i,\qquad b=\ell_k,\qquad c=\ell_j=L-a-b,
 \qquad n=\min(a,c)
\]
と書く．固定した $\sigma>0$ と有界区間 $J^+$ に対し，残余 collar を
\[
 q=\log(b/b_0)\in J^+,\qquad \theta=\arg\zeta_k,
 \qquad b_0\longrightarrow\infty
\]
で表す．外側領域では $a,c>(1+\sigma)b$ とし，全 phase を含める．
Normal subgroup を weights $(1,-1,0)$ の $\alpha$，その補完を
weights $(0,-1,1)$ の $\beta$ とする．
共通座標に移した actual form と小さい canonical primitive を $\omega=d\lambda$ とし，
\[
 \eta=\mathcal A_\alpha\omega,\qquad
 \kappa=\mathcal A_\beta\eta,\qquad
 \mu=(\mathcal A_\alpha\lambda)(V_\alpha),\qquad
 \nu=\mathcal A_\beta\mu
\]
と置く．$\psi$ は centered homotopy による第二段の map，すなわち
$\psi^*\eta=\kappa$，$\mu\circ\psi=\nu$ を満たすものとする．

共通 interval $I$ が与えられ，各 $u\in I$ と各 $q\in J^+$ について
$\nu=u$ は normal radial line 上に唯一の level を持ち，その level は
$a,c>(1+2\sigma)b$ の内側に留まると仮定する．
残余 collar にも固定の正の $q$-幅の余裕を取り，必要な比較が成立する
外側 arm に全ての level とその一様な $h=g_L^{\rm sf}$ 近傍が含まれるものとする．
Interval $I$ の存在はこの命題の仮定である．

このとき，大きい $b_0$ について $\eta$ の対応する reductions も
残余座標 $(q,\theta)$ で同定できる．$\psi$ による reduced map は
\begin{equation}\label{ke:eq:reduced-arm-map}
 D_u(q,\theta)=(q+u_0(u,q,\theta),\theta+v_0(u,q,\theta))
\end{equation}
と書け，角変位には恒等写像からの一価な実 lift を取れる．
固定 $u$ での $(q,\theta)$ に関する通常の $C^1$ norm について，$u\in I$ に一様に
\[
 \|u_0(u,\cdot)\|_{C^1}\le Cb_0^{-2},\qquad
 \|v_0(u,\cdot)\|_{C^1}\le Cb_0^{-1}
\]
を満たす．従って $\kappa$ の radial circle $q=s$ の像は
\begin{equation}\label{ke:eq:actual-reduced-circle-graph}
 q=Q_u(s,\theta),\qquad
 \|Q_u(s,\theta)-s\|_{C^1(s,\theta)}\le Cb_0^{-2}
\end{equation}
という graph である．各族は $u$ に滑らかであるが，
ここでの $C^1$ norms に無重みの $\partial_u$ は含めない．
補題\ref{ke:reduced-collar-graph}により，これらの end 葉を固定幅の広い collar 内で
基準 foliation へ相対的に延長できる．
\end{proposition}

\begin{proof}
この固定幅の残余 collar では $b\asymp b_0$ である．
Full-triple 比較と座標接続の指数評価により，$\omega,\eta,\kappa$ は
模型から全ての必要な固定階で $O(b_0^{-1})$ 以内にある．
模型の正規化 chart で最大の normal 深度を従属座標に選ぶと，
$V_\alpha$ の非零係数は残る normal 深度の逆数と比較可能である．
Constrained cusp metric と complete model の一様比較および高階 bounds から，
\begin{equation}\label{ke:eq:normal-weighted-derivatives}
 \mathfrak a:=|V_\alpha|_h\asymp n^{-1},\qquad
 |\nabla_h^rV_\alpha|_h\le C_r\mathfrak a,
 \qquad |\nabla_h^rV_\beta|_h\le C_r b_0^{-1}
\end{equation}
を得る．$\mathfrak a$ の対数微分も全ての必要な固定階で一様有界である．
ここでは $n/b$ の上界を仮定していない．

Centered 第二段 primitive を $\beta_2=K_\beta\eta$ とする．
$K_\beta\omega_0=0$ と\eqref{ke:eq:normal-weighted-derivatives}により
$\|\beta_2\|_{C^3(h)}\le Cb_0^{-2}$，第二段の場と flow の $C^2$ 誤差も
$O(b_0^{-2})$ である．Moment identity の証明で用いた式から
\[
 \beta_2(V_\alpha)=\nu-\mu.
\]
従って $r\le2$ について
\begin{equation}\label{ke:eq:weighted-moment-difference}
 |\nabla_h^r(\mu-\nu)|_h\le C\mathfrak a b_0^{-2}.
\end{equation}
Normal moment の勾配が小さい場合に，右辺の $\mathfrak a$ を落としてはならない．

$N=-J_0V_\alpha/\mathfrak a$ と置く．これは残余座標を保つ単位 normal radial 場で，
その固定階の微分は一様有界である．$\eta,\kappa$ の $J_0$-tame 下界から
\[
 d\mu(N)\ge c_1\mathfrak a,\qquad
 d\nu(N)\ge c_1\mathfrak a,
 \qquad |\nabla_h^r d\mu|_h+|\nabla_h^r d\nu|_h\le C_r\mathfrak a
\]
を得る．$\nu=u$ の点から $N$ に沿って動き，同じ残余座標で $\mu=u$ を解く．
式\eqref{ke:eq:weighted-moment-difference}と上の単調性により，
一意な解までの normal 変位は $Cb_0^{-2}$ 以下である．
正規化した方程式を出発点の $\mathfrak a$ で割れば，normal 微分は一様に正で，
誤差とその固定 $u$ での tangential 微分は $O(b_0^{-2})$ になる．
従って通常の陰関数定理の定量的評価は $n$ に依存しない．
Level の normal 深度が非常に大きくても，$\mathfrak a$ の小ささは分子と分母で相殺される．

使用した tangential bounds を具体的に確認する．
$a$ を固定して $b$ を動かす $Q_0=\partial_q|_a$ は，
$b\le c$ と模型比較から $h$-長さと固定階の微分が一様に有界である．
$\kappa$ level 上の $q$ 方向は
\[
 S=Q_0-\frac{d\nu(Q_0)}{d\nu(N)}N
\]
で与えられる．上の weighted bounds により，その係数と固定階の微分は
一様有界である．従って level の $(q,\theta)$ section において
$|\partial_q|_h\le C$，$|\partial_\theta|_h\le C/b_0$ であり，
必要な微分も同じ尺度で制御される．$\nu$ は全 torus 不変なので，
normal phase を固定してこの section を選べる．
$\eta$ 側では円平均による normal 不変性と先の一意性があり，
同じ残余座標が reduction を parameter 化する．

$\psi$ は normal 円と可換で，$\mu\circ\psi=\nu$ であるから，
この二つの reductions の間に写像を誘導する．
第二段の全 trajectories は固定の arm margin と残余 buffer に留まる．
その場を $Z_s$ とし，$q$ と角の閉一形式 $d\theta$ を trajectory に沿って積分する．
$dq$ の固定階の norm は有界，$d\theta$ の固定階の norm は $O(b_0)$ なので，
\[
 u_0=\int_0^1dq(Z_s)\circ\psi_s\,ds,
 \qquad
 v_0=\int_0^1d\theta(Z_s)\circ\psi_s\,ds
\]
はそれぞれ ambient $C^1(h)$ で $O(b_0^{-2})$，$O(b_0^{-1})$ である．
後者の積分は winding を恒等写像から連続に固定した実 lift である．
上の level section に制限すると，より具体的には
\[
 |\partial_q u_0|\le Cb_0^{-2},\quad
 |\partial_\theta u_0|\le Cb_0^{-3},\qquad
 |\partial_q v_0|\le Cb_0^{-1},\quad
 |\partial_\theta v_0|\le Cb_0^{-2}.
\]
角の周期を無視して ambient 変位を通常の角座標へ移していないことに注意する．

大きい $b_0$ では $\theta\mapsto\theta+v_0(u,s,\theta)$ は degree $1$ の
orientation を保つ円の diffeomorphism である．その逆写像を代入して
$q=s+u_0(u,s,\theta)$ を書き直せば，
\eqref{ke:eq:actual-reduced-circle-graph}を得る．
すべて固定 $u$ の微分評価であり，定数は $u\in I$ に共通である．
一方，無重みの $\partial_u$ には normal 勾配の逆数が現れ得るので，
同じ小ささは主張しない．固定 $t\ne0$ での滑らかな parameter 依存は
非消滅の normal 微分による陰関数定理から従う．

最後に，graph を定義した広い $q$-collar 内で cutoff を選ぶ．
固定幅なら $\|\chi'\|_\infty$ は一様有界であり，
補題\ref{ke:reduced-collar-graph}の条件は大きい $b_0$ で満たされる．
幅 $w_t$ も縮める場合には，少なくとも
$b_0^{-2}(1+w_t^{-1})\to0$ と，trajectory を収める同じ幅の余裕を別に確認する．
\end{proof}

本命題は，共通 $\kappa$-moment strip が既に与えられた単一 arm について，
reduced 葉の小さい graph という入力を actual 解析から検証する．
三つの end と source-inclusive neck を同時に覆う interval と neck 全域の reduction は，
定理\ref{ke:neck-corner-diagonal-matching}と
系\ref{ke:neck-corner-saturated-union}で別に同定する．
全曲面の接着と距離比較には，さらに定理\ref{ke:actual-global-fibration}と
定理\ref{ke:global-ambient-lower}，\ref{ke:global-ambient-upper}を用いる．
独立に先指定した第二 period と critical actions の数値の一致は，これらの結論にも含めない．

\begin{proposition}[単一 arm の共通 canonical moment strip]
\label{ke:arm-common-moment-strip}
命題\ref{ke:single-arm-reduced-graph}の座標と normal 作用を用い，
$b\asymp b_0$ の残余 annulus に対して
\[
 b_0\longrightarrow\infty,\qquad b_0/L\longrightarrow0
\]
とする．固定定数 $K>1+2\sigma$ を取り，normal radial interval を
\[
 Kb\le a\le L-(K+1)b
 \qquad\bigl(c=L-a-b\ge Kb\bigr)
\]
とする．全ての座標は内側の正則 plateau arm にあり，
triple からの小さい primitive を共通に用いる．このとき両端で，
全角度および残余 annulus に一様に
\begin{align}
 \mu|_{a=Kb},\ \nu|_{a=Kb}
  &=\frac{4\pi}{Kb}
       +O\!\left(b_0^{-2}+\frac{(b_0/L)^2}{b_0}\right),
       \label{ke:eq:arm-positive-moment}\\
 \mu|_{c=Kb},\ \nu|_{c=Kb}
  &=-\frac{4\pi}{Kb}
       +O\!\left(b_0^{-2}+\frac{(b_0/L)^2}{b_0}\right).
       \label{ke:eq:arm-negative-moment}
\end{align}
従って，ある $c_*>0$ に対し
\[
 I_t=[-c_*/b_0,c_*/b_0]
\]
の全ての値は，各残余点において $\nu$ と $\mu$ のそれぞれ唯一の normal level を持つ．
それらは上の interval の内部にあり，dominant arm の固定比の margin を持つ．
残余 annulus に固定幅の余裕を取れば，
命題\ref{ke:single-arm-reduced-graph}の共通 strip の仮定をこの $I_t$ で満たす．
従ってこの単一 arm では，共通 interval と actual reduced graph の両方が得られる．
\end{proposition}

\begin{proof}
$a=Kb$ の端では $(a+b)/L=(K+1)b/L\to0$ であり，$c$ が最大深度である．
まず座標接続前の共通 triple chart で比較する．
独立座標を $(\zeta_i,\zeta_k)$ とし，対応する二つの cusp の積を
$\omega_{\rm prod}$ とする．関連 KE 主稿\cite[
\texttt{simplex:prop:dominant-product}]{KELocalCurrent}と，
同稿の系 \texttt{simplex:cor:actual-dominant-product} の
固定半径正規化 chart 上の全固定階評価により，
\[
 \left\|\log\frac{\omega_t^2}{\omega_{\rm prod}^2}\right\|_{C^2}
 \le C\left(b_0^{-1}+(b_0/L)^2\right)
\]
を得る．射影座標へ移した場合も，命題\ref{ke:three-arm-coordinate-interpolation}の
指数的な $C^q$ 誤差を加えるだけであり，それは $Cb_0^{-1}$ に吸収できる．
ここで比較しているのは canonical volume potentials であり，
任意の primitive の差を tensor 比較から結論していない．

新しい logarithmic frame $\Omega_{\log}=d\log\zeta_i\wedge d\log\zeta_k$ に対して
\[
 \lambda_{\rm prod}
 =d^c\log\frac{\omega_{\rm prod}^2}{|\Omega_{\log}|^2}
 =d^c(-2\log a-2\log b)
\]
である．正の定数による volume の規格化は $d^c$ で消える．
周期 $1$ の $V_\alpha$ は独立座標の $\arg\zeta_i$ を $2\pi$ だけ回すから
\[
 \lambda_{\rm prod}(V_\alpha)=\frac{4\pi}{a}.
\]
上の volume-potential 差に $d^c$ を施し，
$|V_\alpha|=O(b_0^{-1})$ を掛ければ，normal moment の誤差は
\eqref{ke:eq:arm-positive-moment}の率になる．

座標変更後の小さい primitive が，新座標の canonical primitive と
厳密には異なる点も確認する．正則 plateau で
$G^*\Omega_{\rm old}=J_G\Omega_{\log}$ と書くと，
$J_G=1+O(e^{-c b_0})$ は一価な正則対数を持ち，二つの primitive の差は
$-d\operatorname{Im}\operatorname{Log}J_G$ である．
この exact 項の normal 円上の平均積分は零である．
また積模型の表示は全 torus 不変なので，normal 平均，さらに残余平均を取っても
同じ leading term と誤差が得られる．従って $\mu$，$\nu$ の双方に上の式が成立する．
射影 neck の canonical moment との比較には
系\ref{ke:transition-normal-primitive}を使い，加法定数を独立に選び直さない．

$c=Kb$ の端では独立座標を $(\zeta_j,\zeta_k)$ に取り直す．
Logarithmic frame は符号だけ変わり，absolute volume は同じである．
$V_\alpha$ の $\arg\zeta_j$ 成分は $-2\pi$ なので，
同じ計算から\eqref{ke:eq:arm-negative-moment}を得る．

残余 collar 全体で $b\le C_+b_0$ とする．例えば
$c_*=\pi/(KC_+)$ と固定すれば，十分大きい $b_0$，十分小さい $b_0/L$ に対し，
正の端の moment は $3c_*/b_0$ 以上，負の端は $-3c_*/b_0$ 以下となる．
このとき $L>(2K+1)b$ なので normal interval は空でない．
前命題の $J_0$-tame 評価により，$\mu$，$\nu$ はこの normal radial line に沿って
狭義単調である．従って中間値の定理から $I_t$ の各値の存在と一意性を得る．
全 level は $a,c\ge Kb>(1+2\sigma)b$ を満たし，
外側 $a,c>(1+\sigma)b$ から固定比の余裕がある．
少し広い残余 annulus でも同じ評価を先に適用し，内側に制限すれば，
残余方向の正の $h$-buffer も取れる．これで前命題の適用条件を確認した．
\end{proof}

この interval は三重点の一つの source-free arm に関するものである．
Source-inclusive neck 側の interval と共通の開領域を選び，
両側の normal trajectories をそこに収めることは，
定理\ref{ke:neck-corner-diagonal-matching}の exhaustion と buffer の評価によって行う．
ここでの単一 arm の証明に exterior の $(\mathrm{TC})_d$ や
global $(\mathrm{CP})_d+(\mathrm{PBE})_d$ は用いていない．

\subsection{座標同定に関する normal Moser flow の自然性}

\begin{lemma}[共通の円作用に対する chain homotopy と flow の自然性]
\label{ke:normal-flow-naturality}
Physical neck 上の円作用を $\rho_N$，その閉 symplectic form を $\omega_N$ とする．
新しい座標領域 $W$ は標準円作用 $\rho_\alpha$ に飽和しているとし，
embedding $G:W\to N$ が，全ての $z\in W$ と $\theta\in\RR/\ZZ$ について
\[
 G\rho_{\alpha,\theta}(z)=\rho_{N,\theta}G(z)
\]
を満たすとする．$\omega_\alpha=G^*\omega_N$ と置き，両側で同じ
centered kernel $1/2-\theta$ による $K_N,K_\alpha$ と円平均
$\mathcal A_N,\mathcal A_\alpha$ を使う．このとき
\begin{equation}\label{ke:eq:normal-homotopy-naturality}
 G^*\mathcal A_N\omega_N=\mathcal A_\alpha\omega_\alpha,
 \qquad G^*K_N\omega_N=K_\alpha\omega_\alpha.
\end{equation}
両側の normal Moser path
$\Omega_{N,s}=(1-s)\mathcal A_N\omega_N+s\omega_N$ と
$\Omega_{\alpha,s}=(1-s)\mathcal A_\alpha\omega_\alpha+s\omega_\alpha$
が非退化なら，対応する Moser vector fields は $G$ で関係付けられる．
従って共通の定義域からの全 trajectories が $W$ と $G(W)$ に留まる限り，
\begin{equation}\label{ke:eq:normal-flow-conjugacy}
 \phi_{N,s}\circ G=G\circ\phi_{\alpha,s}
 \qquad(0\le s\le1)
\end{equation}
が成立する．この flow の一致には，二つの canonical primitive の一致を要しない．
\end{lemma}
\begin{proof}
Equivariance を微分すると $dG(V_\alpha)=V_N\circ G$ である．
従って任意の形式 $\xi$ に対して
\[
 G^*\rho_{N,\theta}^*\xi
 =\rho_{\alpha,\theta}^*G^*\xi,\qquad
 G^*\iota_{V_N}\xi=\iota_{V_\alpha}G^*\xi.
\]
同じ kernel を掛けて全円上で積分すれば
\eqref{ke:eq:normal-homotopy-naturality}を得る．
$G^*\Omega_{N,s}=\Omega_{\alpha,s}$ と非退化性により，
$\iota_{Y_s}\Omega_s=K\omega$ の解も $G$ で関係付けられる．
初期値は両側で $G(z)$ なので，ODE の一意性から
\eqref{ke:eq:normal-flow-conjugacy}が従う．
使っている primitive は形式 $\omega$ から直接作る $K\omega$ であり，
$\omega=d\lambda$ の $\lambda$ の選択ではない．
また，平均には全 circle orbit の情報が必要である．
一点での作用生成場の一致や，共通領域から退出する trajectory に
ODE の一致を適用することはしていない．
\end{proof}

本稿への適用では，$F$ を命題\ref{ke:three-arm-coordinate-interpolation}の
写像，$G=F^{-1}$ とする．その内側 plateau における
$\rho_N$ は，補題\ref{ke:projected-source-normal-form}の
$(x',y')$ を逆向きに回す source-inclusive projected action である．
全 circle orbit を含み，さらに二つの normal flow と必要な residual stage の
像を収める開集合を plateau 内に取れば，上の補題をそのまま適用できる．
固定した深度比と残余幅の余裕の下では，
命題\ref{ke:quantitative-order-hamiltonian}と
命題\ref{ke:single-arm-reduced-graph}の小変位が局所の滞在条件を与える．
その開集合が global neck の定義域にも含まれることは，別に確認する．

特に source-inclusive neck の構成にも同じ centered kernel を選ぶ．
通常 kernel による既存の endpoint と同一だとは主張せず，許される構成の選択を
両側でそろえるのである．Global bundle frame から局所 frame への変更は
base-only であり，normal circle は変わらない．
体積 frame の変更も残余基底の非零因子によるため，normal moment の規格化は変わらない．

\begin{corollary}[同じ actual normal action 上の corner 葉]
\label{ke:corner-leaves-in-neck-reduction}
前補題の共通開領域で，$\eta_N=\mathcal A_N\omega_N$，
$\eta_\alpha=\mathcal A_\alpha\omega_\alpha$ とする．
Neck の商が共通 interval $I$ 上で $(\mu_N,p)$ によって
$I\times P$ と同定されるとする．ここで使う end は source を含まず，
その reduction は滑らかである．Corner 側には小さい primitive
$\lambda_\alpha=G^*\lambda_\zeta$ を用い，
$\mu_\alpha=(\mathcal A_\alpha\lambda_\alpha)(V_\alpha)$ とする．
第二段を $\psi$ とし，
\[
 \psi^*\eta_\alpha=\kappa,\qquad
 \mu_\alpha\circ\psi=\nu
\]
とする．全ての残余 stage の像およびその後の normal trajectories が
前補題の共通領域に留まると仮定する．
Canonical normal moment を系\ref{ke:transition-normal-primitive}で規格化すれば，
\[
 \mu_\alpha=\mu_N\circ G,\qquad
 G\phi_\alpha\psi=\phi_NG\psi.
\]
従って $\nu=u$ 上の標準 torus $T$ に対し，$G\psi(T)$ は
$\eta_N$ の同じ normal reduction 内の残余円
$\gamma_{u,T}=p(G\psi(T))$ の全 normal-circle 逆像であり，
actual corner torus は厳密に
\begin{equation}\label{ke:eq:corner-neck-same-fiber}
 G\phi_\alpha\psi(T)
 =\phi_N\bigl((\mu_N,p)^{-1}(\{u\}\times\gamma_{u,T})\bigr)
\end{equation}
である．右辺の逆像は同じ neck moment strip 内で取る．
\end{corollary}
\begin{proof}
Normal moment の等式は系\ref{ke:transition-normal-primitive}による．
Global frame の base-only 因子も normal 円上の積分を変えない．
Flow の共役式は前補題を $\psi(T)$ の各点に適用して得られる．
$\psi$ は $\rho_\alpha$ と可換であり，$G$ はこの円を $\rho_N$ へ移すので，
$G\psi(T)$ は全 normal-circle orbit で飽和している．
また $\mu_N\circ G\psi=\nu$ なので同じ moment level に入る．
Reduction の各 $(u,p)$ 上には一つの normal orbit があるため，
残余円の全逆像との等式が成立する．
これに同じ $\phi_N$ を適用すれば\eqref{ke:eq:corner-neck-same-fiber}となる．
\end{proof}

この残余円を読む座標は，元の $p$ の任意の radial coordinate ではない．
新座標の残余成分は
$Z_k=\Xi_{kk}(p)=C_D(p,t)^2/h_t(p)$ なので，
命題\ref{ke:single-arm-reduced-graph}の座標は
\[
 q=\log\!\left(\frac{-\log|\Xi_{kk}(p)|^2}{b_0}\right),
 \qquad \theta=\arg\Xi_{kk}(p)
\]
である．この局所正則座標で得た graph を neck の残余曲面へ戻す．

\subsection{Moment に依存する残余 foliation の許容}

\begin{lemma}[Collar での parameter 付き foliation の変更]
\label{ke:moment-dependent-neck-foliation}
定理\ref{ke:projected-source-neck}の円平均後の領域と共通 moment interval $I$ を取り，
$(\mu_N,p)$ による商を $I\times P$ と同定する．
Source のある $u=0$ での商同定は位相的なものであり，
全 reduction が滑らかであるとは仮定しない．
$\mathcal F_0$ を同定理の marked Morse foliation とする．
固定 $t\ne0$ で $u\in I$ に滑らかに依存する diffeomorphisms
$D_u\in\operatorname{Diff}_0(P)$ が与えられ，
source と saddle の固定近傍，および指定した end collars の外では恒等とする．
このとき $\mathcal F_u=(D_u)_*\mathcal F_0$ の葉を各 $\mu_N=u$ で
持ち上げ，同じ normal Moser map $\phi_N$ で移すと，一つの actual
Lagrangian neck fibration を得る．Regular torus，saddle fibre，
一重 pinched torus の型は元の定理と同じである．
退化 parameter に一様な無重みの $\partial_uD_u$ の評価は要求しない．

特に命題\ref{ke:single-arm-reduced-graph}の graph が共通 $I$ と
開 end collars 上に与えられれば，補題\ref{ke:reduced-collar-graph}の
$D_u$ を用いて，系\ref{ke:corner-leaves-in-neck-reduction}の
corner tori と同じ fibres を開 overlap 全体で得られる．
\end{lemma}
\begin{proof}
$\mathcal F_0$ の葉空間への写像を $f_0:P\to T$ と書けば，
円平均後の neck 写像は
\[
 x\longmapsto
 \bigl(\mu_N(x),f_0(D_{\mu_N(x)}^{-1}(p(x)))\bigr)
\]
で定義される．これは連続で，regular 領域では滑らかである．
各 fibre では $\mu_N=u$ を固定しているので，その接空間は
$V_N$ と $\ker d\mu_N$ 内の葉方向で張られる．
$\iota_{V_N}\eta_N=-d\mu_N$ により二次元の regular fibre は Lagrangian である．
Parameter に関する微分は fibre の接空間には入らない．
Regular 領域での写像の rank は，$d\mu_N$ と残余曲面上の
非零の葉横断微分から $2$ である．

各固定 $u$ で $D_u$ は diffeomorphism なので葉の連結性と marked graph を保つ．
Source と saddle の近傍は全 $u$ で固定しているため，
元の focus--focus germ と saddle の局所模型を変更しない．
各 source は対応する一つの regular leaf 上に留まり，他の source と混ざらない．
従って全特異 fibre の位相型と各 smooth stratum の isotropic 性も
定理\ref{ke:projected-source-neck}と同じ証明で得られる．
同じ $\phi_N$ による移送で actual form に戻す．

最後の適用では，固定 $u$ の $C^1$ graph 誤差が小さいことから
補題\ref{ke:reduced-collar-graph}が各 $D_u$ を構成する．
Graph の $u$ 依存が滑らかなら $D_u$ も滑らかであり，
collars の外を恒等にするので source と saddle の条件は保たれる．
その開 collar 上で $\mathcal F_u$ の葉は $\gamma_{u,T}$ そのものとなるので，
\eqref{ke:eq:corner-neck-same-fiber}により fibre 分割が厳密に一致する．
ここでも指定された第二 period や critical action の値への独立な調整を
自動的に行ったとは主張しない．
\end{proof}
